\chapter{Functional Programming, Comprehensions \& Generators}

Python supports \textbf{Functional Programming} idioms. Instead of explicitly writing multi-line \texttt{for} loops, functional Python allows developers to write concise, expressive expressions using list comprehensions, anonymous \texttt{lambda} functions, higher-order utilities (\texttt{map}, \texttt{filter}, \texttt{zip}), and memory-efficient \texttt{yield} generators.

\section{List, Dictionary, and Set Comprehensions}

Comprehensions provide a compact 1-line syntax to construct new sequences from existing iterables.

\begin{formulabox}{Comprehension Syntax}
$$\text{List Comprehension: } [\text{expression } \mathbf{for} \text{ item } \mathbf{in} \text{ iterable } \mathbf{if} \text{ condition}]$$
$$\text{Dict Comprehension: } \{\text{key\_expr}: \text{val\_expr } \mathbf{for} \text{ item } \mathbf{in} \text{ iterable}\}$$
\end{formulabox}

\begin{examplebox}{Filtering and Transforming}
\begin{lstlisting}
numbers = [1, 2, 3, 4, 5, 6, 7, 8, 9, 10]

# List Comprehension: Squares of even numbers
even_squares = [x**2 for x in numbers if x % 2 == 0]
print(even_squares)  # Output: [4, 16, 36, 64, 100]

# Dict Comprehension: Number to its cube
cubes = {x: x**3 for x in range(1, 6)}
print(cubes)  # Output: {1: 1, 2: 8, 3: 27, 4: 64, 5: 125}
\end{lstlisting}
\end{examplebox}

\begin{videocard}{IITM BS Lecture: Inline Statements and List Comprehensions}{https://www.youtube.com/watch?v=MNe03MfOLto}
Official lecture explaining syntax, conditional filtering, and nested list comprehensions.
\end{videocard}

\section{Lambda Functions, \texttt{map()}, \texttt{filter()}, \texttt{zip()}, and \texttt{enumerate()}}

\begin{definitionbox}{Lambda and Higher-Order Functions}
\begin{itemize}
    \item \textbf{Lambda (\texttt{lambda x: expr}):} Small anonymous 1-line inline functions.
    \item \textbf{\texttt{map(func, iterable}):} Applies \texttt{func} to every item in \texttt{iterable}.
    \item \textbf{\texttt{filter(predicate, iterable}):} Retains items where \texttt{predicate(item)} is \texttt{True}.
    \item \textbf{\texttt{zip(seq1, seq2}):} Pairs elements from multiple sequences element-wise into tuples.
    \item \textbf{\texttt{enumerate(seq}):} Yields \texttt{(index, item)} pairs during iteration.
\end{itemize}
\end{definitionbox}

\begin{examplebox}{Functional Pipeline Example}
\begin{lstlisting}
nums = [1, 2, 3, 4, 5]

# Square all numbers using map and lambda
squares = list(map(lambda x: x**2, nums))  # [1, 4, 9, 16, 25]

# Filter odd numbers
odds = list(filter(lambda x: x % 2 != 0, nums))  # [1, 3, 5]

# Zip two lists together
names = ["Alice", "Bob"]
scores = [85, 92]
paired = list(zip(names, scores))  # [('Alice', 85), ('Bob', 92)]
\end{lstlisting}
\end{examplebox}

\textbf{Data Science Relevance:} Feature extraction pipelines (like \texttt{df['col'].apply(lambda x: clean(x))}) and model training loops (\texttt{for idx, (batch\_x, batch\_y) in enumerate(dataloader)}) heavily utilize these exact tools.

\begin{videocard}{IITM BS Lecture: Lambda, Enumerate, Zip, Map, and Filter}{https://www.youtube.com/watch?v=ueIignqReDY}
Official lecture covering functional iteration helpers and lambda syntax.
\end{videocard}

\section{Generators and the \texttt{yield} Keyword}

Unlike regular functions that return a full dataset at once (storing everything in RAM), \textbf{Generators} compute items lazily on demand using the \texttt{yield} keyword.

\begin{theorembox}{Generators vs Functions}
A function with \texttt{yield} pauses execution after sending a value back and resumes right where it left off when the next item is requested via \texttt{next()} or a \texttt{for} loop.
\end{theorembox}

\begin{examplebox}{Infinite Stream Generator}
\begin{lstlisting}
def fibonacci_stream():
    a, b = 0, 1
    while True:
        yield a
        a, b = b, a + b

fib = fibonacci_stream()
for _ in range(6):
    print(next(fib), end=" ")  # Output: 0 1 1 2 3 5
print()
\end{lstlisting}
\end{examplebox}

\begin{videocard}{IITM BS Lecture: Iterators, Generators, and Yield}{https://www.youtube.com/watch?v=29A7HhnzxPU}
Official lecture explaining memory optimization, lazy evaluation, and \texttt{yield} semantics.
\end{videocard}

\newpage
\section{Practice Exercises}
\textit{Detailed step-by-step solutions are available in the Appendix.}

\begin{enumerate}
    \item \textbf{List Comprehension Filtering:} Given \texttt{words = ["python", "data", "science", "ai", "machine", "ml"]}, write a 1-line list comprehension to extract words that have a length greater than $3$ characters, converted to UPPERCASE.

    \item \textbf{Lambda Sorting:} Given a list of tuples representing student names and test scores \texttt{students = [("Alice", 88), ("Bob", 95), ("Charlie", 78), ("David", 90)]}, sort the list in descending order based on their test scores using \texttt{sorted()} and a \texttt{lambda} key.

    \item \textbf{Zipping Vectors:} Given two lists \texttt{x = [1, 2, 3]} and \texttt{y = [10, 20, 30]}, write a 1-line list comprehension using \texttt{zip(x, y)} to compute the element-wise sum \texttt{[11, 22, 33]}.

    \item \textbf{Fibonacci Generator:} Write a generator function \texttt{generate\_squares(n)} that yields the squares of numbers from $1$ up to $n$ one by one. Iterate over it using a \texttt{for} loop.

    \item \textbf{Data Science Application (Text Preprocessing Pipeline):} In Natural Language Processing (NLP), a raw text string \texttt{text = "   Data Science and Machine Learning 101   "} needs cleaning. Write a functional pipeline using \texttt{map()}, \texttt{filter()}, and string methods to:
    \begin{enumerate}
        \item Strip whitespace and convert to lowercase.
        \item Tokenize into words.
        \item Filter out words that are numbers (using \texttt{.isdigit()}).
    \end{enumerate}
\end{enumerate}